\section{Metodología}

\subsection{Línea de investigación}
Desarrollo y aplicaciones de la inteligencia artificial.

\subsection{Enfoque, tipo y nivel}
Enfoque cuantitativo; investigación aplicada o de desarrollo tecnológico; nivel
explicativo y comparativo.

\subsection{Diseño}
Diseño experimental comparativo para la fase técnica, con manipulación de la
configuración del detector y baseline YOLOv8n FP32, y evaluación pedagógica
cuasiexperimental de medidas repetidas para la fase docente. La asignación
aleatoria se aplica a las particiones y semillas de los datos, no a los sujetos.

\subsection{Técnicas e instrumentos de recolección}
Para la fase técnica se usan scripts de entrenamiento y evaluación, con registro
de métricas y del hardware. Para la fase pedagógica se emplean fichas
cronométricas, la escala NASA-TLX \parencite{hart1988} y el cuestionario SUS
\parencite{brooke1996}.

\subsection{Procesamiento y análisis}
Se reportan medias y desviaciones entre pliegues. Se aplican Wilcoxon, McNemar,
bootstrap de 5000 remuestreos y un criterio de no inferioridad para H2. Se
documentan versiones de Python, PyTorch, Ultralytics, ONNX Runtime y Android,
además de las semillas y los pesos.
